\documentclass[11pt,a4paper]{article}

\newcommand{\doctitle}{GEV Hazard Curve Model (MKM-GH-001)}
\newcommand{\docsubtitle}{Wolfram Oracle Benchmark of the GEV Implementation}
\newcommand{\docversion}{0.1}
\newcommand{\docdate}{19-October-2026}
\newcommand{\docauthor}{David K Kelly}

% MKM Research Labs — Shared Document Header
% Usage: % MKM Research Labs — Shared Document Header
% Usage: % MKM Research Labs — Shared Document Header
% Usage: \input{../shared/mkm_header}  (before \begin{document})
%
% Required commands before \input:
%   \newcommand{\doctitle}{...}
%   \newcommand{\docsubtitle}{...}
%   \newcommand{\docversion}{1.0}
%   \newcommand{\docdate}{13-February-2026}
%   \newcommand{\docauthor}{David K Kelly}

% ─── Packages ───────────────────────────────────────────────────────────────
\usepackage[utf8]{inputenc}
\usepackage{amsmath,amsfonts,amssymb}
\usepackage{listings}
\usepackage{color,xcolor}
\usepackage{graphicx}
\usepackage{hyperref}
\usepackage[toc,page]{appendix}
\usepackage{tabto}
\usepackage{float}
\usepackage{comment}
\usepackage[en-GB]{datetime2}
\usepackage{setspace}
\usepackage{subcaption}
\usepackage{booktabs}
\usepackage{longtable}
\usepackage{adjustbox}
\usepackage{enumitem}
\usepackage{fancyhdr}
\usepackage{lastpage}

% ─── Page geometry ──────────────────────────────────────────────────────────
\onehalfspacing
\usepackage{geometry}
\geometry{a4paper, total={170mm,247mm}, left=20mm, top=25mm, bottom=25mm}

% ─── Headers and footers ───────────────────────────────────────────────────
\pagestyle{fancy}
\fancyhf{}
\fancyhead[L]{\small\textcolor{gray}{MKM Research Labs}}
\fancyhead[R]{\small\textcolor{gray}{\doctitle{} --- v\docversion}}
\fancyfoot[C]{\small\textcolor{gray}{CONFIDENTIAL}}
\fancyfoot[R]{\small\textcolor{gray}{Page \thepage\ of \pageref{LastPage}}}
\renewcommand{\headrulewidth}{0.4pt}
\renewcommand{\footrulewidth}{0.4pt}

% ─── Colours ────────────────────────────────────────────────────────────────
\definecolor{mkmblue}{HTML}{1565C0}
\definecolor{codegreen}{rgb}{0,0.6,0}
\definecolor{codegray}{rgb}{0.5,0.5,0.5}
\definecolor{codepurple}{rgb}{0.58,0,0.82}
\definecolor{backcolour}{rgb}{0.95,0.95,0.92}

% ─── Hyperref setup ─────────────────────────────────────────────────────────
\hypersetup{
    colorlinks=true,
    linkcolor=mkmblue,
    citecolor=mkmblue,
    urlcolor=mkmblue
}

% ─── Code listings ──────────────────────────────────────────────────────────
\lstdefinestyle{mystyle}{
    backgroundcolor=\color{backcolour},
    commentstyle=\color{codegreen},
    keywordstyle=\color{magenta},
    numberstyle=\tiny\color{codegray},
    stringstyle=\color{codepurple},
    basicstyle=\ttfamily\footnotesize,
    breakatwhitespace=false,
    breaklines=true,
    captionpos=b,
    keepspaces=true,
    numbers=left,
    numbersep=5pt,
    showspaces=false,
    showstringspaces=false,
    showtabs=false,
    tabsize=2
}
\lstset{style=mystyle}

% ─── Section formatting ────────────────────────────────────────────────────
\setcounter{tocdepth}{3}
\setcounter{secnumdepth}{3}

% ─── Shared macros ──────────────────────────────────────────────────────────
\newcommand{\mkmtitlepage}{%
    \begin{titlepage}
    \begin{center}
    \vspace*{1cm}
    \begin{figure}[H]\centering
        \includegraphics[width=0.4\textwidth]{../shared/Logos/MKM_logo.png}
    \end{figure}
    \vspace{1.5cm}
    {\Huge\bfseries\textcolor{mkmblue}{\doctitle}\par}
    \vspace{0.5cm}
    {\Large\textcolor{gray}{\docsubtitle}\par}
    \vspace{2cm}
    {\large MKM Research Labs\par}
    \vspace{0.5cm}
    {\large Version \docversion{} --- \docdate\par}
    \vspace{0.5cm}
    {\large \docauthor\par}
    \vspace{2cm}
    \textcolor{red}{\textbf{CONFIDENTIAL --- PROPRIETARY}}
    \end{center}
    \end{titlepage}
}

\newcommand{\mkmlegalpage}{%
    \clearpage
    \thispagestyle{empty}
    \begin{center}
    \large\textbf{Legal Notice}
    \vspace{1em}

    \textbf{PROPRIETARY AND CONFIDENTIAL}
    \end{center}

    \noindent This document, including all algorithms, models, methodology, formulas,
    calculations, and intellectual concepts contained herein, constitutes the exclusive
    intellectual property and confidential trade secrets of MKM Research Labs (``MKM'').

    \vspace{0.5em}
    \noindent Any usage, reproduction, distribution, modification, or disclosure of this
    document or any portion thereof, without the express prior written authorisation from
    MKM Research Labs is strictly prohibited and constitutes an infringement of intellectual
    property rights.

    \vspace{0.5em}
    \noindent The document is provided on an ``as is'' basis. MKM Research Labs makes no
    representations or warranties of any kind, express or implied, regarding its accuracy,
    reliability, or suitability for any particular purpose.

    \vspace{1em}
    \begin{center}
    All rights reserved. \copyright\ 2021--2026 MKM Research Labs.

    \vspace{0.5em}
    Governed by the laws of the United Kingdom.
    \end{center}
    \clearpage
}

\newcommand{\mkmhistorypage}[1]{%
    \section{Document History}
    \begin{table}[ht]
        \centering
        \begin{tabular}{l|l|l|l}
            \textbf{Date} & \textbf{Version} & \textbf{Description} & \textbf{Author} \\
            \hline
            #1
        \end{tabular}
        \label{tab:revision_history}
    \end{table}
    \clearpage
}

% ─── Bank validation protocol section commands ──────────────────────────────
% Use these to ensure every model doc has the required sections.
% They are just styled \section calls — content is provided by each doc.

\newcommand{\mkmexecsummary}{\section{Executive Summary}}
\newcommand{\mkmpurpose}{\section{Model Purpose and Scope}}
\newcommand{\mkmframework}{\section{Mathematical Framework}}
\newcommand{\mkminputs}{\section{Input Parameters and Calibration}}
\newcommand{\mkmimplementation}{\section{Implementation Details}}
\newcommand{\mkmsensitivity}{\section{Sensitivity Analysis}}
\newcommand{\mkmvalidation}{\section{Validation and Backtesting}}
\newcommand{\mkmtestresults}{\subsection{Automated Test Results}\input{test_results}}
\newcommand{\mkmlimitations}{\section{Model Limitations and Known Weaknesses}}

\newcommand{\mkmgovernance}{\section{Governance Validation Questions}
\noindent The following nine questions are mandated by Chapter~5 of the
\emph{Model Risk Governance for Vendors} handbook and must be explicitly
addressed for every model in the inventory. Response status is tracked in the
Model Governance Dashboard (MRC portal). \par\medskip}

\newcommand{\mkmvqitem}[3]{\subsection{Q#1 --- #2}
\noindent\textit{``#3''}\par\medskip\noindent}

\newcommand{\mkmchangelog}{\section{Change History}}
  (before \begin{document})
%
% Required commands before \input:
%   \newcommand{\doctitle}{...}
%   \newcommand{\docsubtitle}{...}
%   \newcommand{\docversion}{1.0}
%   \newcommand{\docdate}{13-February-2026}
%   \newcommand{\docauthor}{David K Kelly}

% ─── Packages ───────────────────────────────────────────────────────────────
\usepackage[utf8]{inputenc}
\usepackage{amsmath,amsfonts,amssymb}
\usepackage{listings}
\usepackage{color,xcolor}
\usepackage{graphicx}
\usepackage{hyperref}
\usepackage[toc,page]{appendix}
\usepackage{tabto}
\usepackage{float}
\usepackage{comment}
\usepackage[en-GB]{datetime2}
\usepackage{setspace}
\usepackage{subcaption}
\usepackage{booktabs}
\usepackage{longtable}
\usepackage{adjustbox}
\usepackage{enumitem}
\usepackage{fancyhdr}
\usepackage{lastpage}

% ─── Page geometry ──────────────────────────────────────────────────────────
\onehalfspacing
\usepackage{geometry}
\geometry{a4paper, total={170mm,247mm}, left=20mm, top=25mm, bottom=25mm}

% ─── Headers and footers ───────────────────────────────────────────────────
\pagestyle{fancy}
\fancyhf{}
\fancyhead[L]{\small\textcolor{gray}{MKM Research Labs}}
\fancyhead[R]{\small\textcolor{gray}{\doctitle{} --- v\docversion}}
\fancyfoot[C]{\small\textcolor{gray}{CONFIDENTIAL}}
\fancyfoot[R]{\small\textcolor{gray}{Page \thepage\ of \pageref{LastPage}}}
\renewcommand{\headrulewidth}{0.4pt}
\renewcommand{\footrulewidth}{0.4pt}

% ─── Colours ────────────────────────────────────────────────────────────────
\definecolor{mkmblue}{HTML}{1565C0}
\definecolor{codegreen}{rgb}{0,0.6,0}
\definecolor{codegray}{rgb}{0.5,0.5,0.5}
\definecolor{codepurple}{rgb}{0.58,0,0.82}
\definecolor{backcolour}{rgb}{0.95,0.95,0.92}

% ─── Hyperref setup ─────────────────────────────────────────────────────────
\hypersetup{
    colorlinks=true,
    linkcolor=mkmblue,
    citecolor=mkmblue,
    urlcolor=mkmblue
}

% ─── Code listings ──────────────────────────────────────────────────────────
\lstdefinestyle{mystyle}{
    backgroundcolor=\color{backcolour},
    commentstyle=\color{codegreen},
    keywordstyle=\color{magenta},
    numberstyle=\tiny\color{codegray},
    stringstyle=\color{codepurple},
    basicstyle=\ttfamily\footnotesize,
    breakatwhitespace=false,
    breaklines=true,
    captionpos=b,
    keepspaces=true,
    numbers=left,
    numbersep=5pt,
    showspaces=false,
    showstringspaces=false,
    showtabs=false,
    tabsize=2
}
\lstset{style=mystyle}

% ─── Section formatting ────────────────────────────────────────────────────
\setcounter{tocdepth}{3}
\setcounter{secnumdepth}{3}

% ─── Shared macros ──────────────────────────────────────────────────────────
\newcommand{\mkmtitlepage}{%
    \begin{titlepage}
    \begin{center}
    \vspace*{1cm}
    \begin{figure}[H]\centering
        \includegraphics[width=0.4\textwidth]{../shared/Logos/MKM_logo.png}
    \end{figure}
    \vspace{1.5cm}
    {\Huge\bfseries\textcolor{mkmblue}{\doctitle}\par}
    \vspace{0.5cm}
    {\Large\textcolor{gray}{\docsubtitle}\par}
    \vspace{2cm}
    {\large MKM Research Labs\par}
    \vspace{0.5cm}
    {\large Version \docversion{} --- \docdate\par}
    \vspace{0.5cm}
    {\large \docauthor\par}
    \vspace{2cm}
    \textcolor{red}{\textbf{CONFIDENTIAL --- PROPRIETARY}}
    \end{center}
    \end{titlepage}
}

\newcommand{\mkmlegalpage}{%
    \clearpage
    \thispagestyle{empty}
    \begin{center}
    \large\textbf{Legal Notice}
    \vspace{1em}

    \textbf{PROPRIETARY AND CONFIDENTIAL}
    \end{center}

    \noindent This document, including all algorithms, models, methodology, formulas,
    calculations, and intellectual concepts contained herein, constitutes the exclusive
    intellectual property and confidential trade secrets of MKM Research Labs (``MKM'').

    \vspace{0.5em}
    \noindent Any usage, reproduction, distribution, modification, or disclosure of this
    document or any portion thereof, without the express prior written authorisation from
    MKM Research Labs is strictly prohibited and constitutes an infringement of intellectual
    property rights.

    \vspace{0.5em}
    \noindent The document is provided on an ``as is'' basis. MKM Research Labs makes no
    representations or warranties of any kind, express or implied, regarding its accuracy,
    reliability, or suitability for any particular purpose.

    \vspace{1em}
    \begin{center}
    All rights reserved. \copyright\ 2021--2026 MKM Research Labs.

    \vspace{0.5em}
    Governed by the laws of the United Kingdom.
    \end{center}
    \clearpage
}

\newcommand{\mkmhistorypage}[1]{%
    \section{Document History}
    \begin{table}[ht]
        \centering
        \begin{tabular}{l|l|l|l}
            \textbf{Date} & \textbf{Version} & \textbf{Description} & \textbf{Author} \\
            \hline
            #1
        \end{tabular}
        \label{tab:revision_history}
    \end{table}
    \clearpage
}

% ─── Bank validation protocol section commands ──────────────────────────────
% Use these to ensure every model doc has the required sections.
% They are just styled \section calls — content is provided by each doc.

\newcommand{\mkmexecsummary}{\section{Executive Summary}}
\newcommand{\mkmpurpose}{\section{Model Purpose and Scope}}
\newcommand{\mkmframework}{\section{Mathematical Framework}}
\newcommand{\mkminputs}{\section{Input Parameters and Calibration}}
\newcommand{\mkmimplementation}{\section{Implementation Details}}
\newcommand{\mkmsensitivity}{\section{Sensitivity Analysis}}
\newcommand{\mkmvalidation}{\section{Validation and Backtesting}}
\newcommand{\mkmtestresults}{\subsection{Automated Test Results}\input{test_results}}
\newcommand{\mkmlimitations}{\section{Model Limitations and Known Weaknesses}}

\newcommand{\mkmgovernance}{\section{Governance Validation Questions}
\noindent The following nine questions are mandated by Chapter~5 of the
\emph{Model Risk Governance for Vendors} handbook and must be explicitly
addressed for every model in the inventory. Response status is tracked in the
Model Governance Dashboard (MRC portal). \par\medskip}

\newcommand{\mkmvqitem}[3]{\subsection{Q#1 --- #2}
\noindent\textit{``#3''}\par\medskip\noindent}

\newcommand{\mkmchangelog}{\section{Change History}}
  (before \begin{document})
%
% Required commands before \input:
%   \newcommand{\doctitle}{...}
%   \newcommand{\docsubtitle}{...}
%   \newcommand{\docversion}{1.0}
%   \newcommand{\docdate}{13-February-2026}
%   \newcommand{\docauthor}{David K Kelly}

% ─── Packages ───────────────────────────────────────────────────────────────
\usepackage[utf8]{inputenc}
\usepackage{amsmath,amsfonts,amssymb}
\usepackage{listings}
\usepackage{color,xcolor}
\usepackage{graphicx}
\usepackage{hyperref}
\usepackage[toc,page]{appendix}
\usepackage{tabto}
\usepackage{float}
\usepackage{comment}
\usepackage[en-GB]{datetime2}
\usepackage{setspace}
\usepackage{subcaption}
\usepackage{booktabs}
\usepackage{longtable}
\usepackage{adjustbox}
\usepackage{enumitem}
\usepackage{fancyhdr}
\usepackage{lastpage}

% ─── Page geometry ──────────────────────────────────────────────────────────
\onehalfspacing
\usepackage{geometry}
\geometry{a4paper, total={170mm,247mm}, left=20mm, top=25mm, bottom=25mm}

% ─── Headers and footers ───────────────────────────────────────────────────
\pagestyle{fancy}
\fancyhf{}
\fancyhead[L]{\small\textcolor{gray}{MKM Research Labs}}
\fancyhead[R]{\small\textcolor{gray}{\doctitle{} --- v\docversion}}
\fancyfoot[C]{\small\textcolor{gray}{CONFIDENTIAL}}
\fancyfoot[R]{\small\textcolor{gray}{Page \thepage\ of \pageref{LastPage}}}
\renewcommand{\headrulewidth}{0.4pt}
\renewcommand{\footrulewidth}{0.4pt}

% ─── Colours ────────────────────────────────────────────────────────────────
\definecolor{mkmblue}{HTML}{1565C0}
\definecolor{codegreen}{rgb}{0,0.6,0}
\definecolor{codegray}{rgb}{0.5,0.5,0.5}
\definecolor{codepurple}{rgb}{0.58,0,0.82}
\definecolor{backcolour}{rgb}{0.95,0.95,0.92}

% ─── Hyperref setup ─────────────────────────────────────────────────────────
\hypersetup{
    colorlinks=true,
    linkcolor=mkmblue,
    citecolor=mkmblue,
    urlcolor=mkmblue
}

% ─── Code listings ──────────────────────────────────────────────────────────
\lstdefinestyle{mystyle}{
    backgroundcolor=\color{backcolour},
    commentstyle=\color{codegreen},
    keywordstyle=\color{magenta},
    numberstyle=\tiny\color{codegray},
    stringstyle=\color{codepurple},
    basicstyle=\ttfamily\footnotesize,
    breakatwhitespace=false,
    breaklines=true,
    captionpos=b,
    keepspaces=true,
    numbers=left,
    numbersep=5pt,
    showspaces=false,
    showstringspaces=false,
    showtabs=false,
    tabsize=2
}
\lstset{style=mystyle}

% ─── Section formatting ────────────────────────────────────────────────────
\setcounter{tocdepth}{3}
\setcounter{secnumdepth}{3}

% ─── Shared macros ──────────────────────────────────────────────────────────
\newcommand{\mkmtitlepage}{%
    \begin{titlepage}
    \begin{center}
    \vspace*{1cm}
    \begin{figure}[H]\centering
        \includegraphics[width=0.4\textwidth]{../shared/Logos/MKM_logo.png}
    \end{figure}
    \vspace{1.5cm}
    {\Huge\bfseries\textcolor{mkmblue}{\doctitle}\par}
    \vspace{0.5cm}
    {\Large\textcolor{gray}{\docsubtitle}\par}
    \vspace{2cm}
    {\large MKM Research Labs\par}
    \vspace{0.5cm}
    {\large Version \docversion{} --- \docdate\par}
    \vspace{0.5cm}
    {\large \docauthor\par}
    \vspace{2cm}
    \textcolor{red}{\textbf{CONFIDENTIAL --- PROPRIETARY}}
    \end{center}
    \end{titlepage}
}

\newcommand{\mkmlegalpage}{%
    \clearpage
    \thispagestyle{empty}
    \begin{center}
    \large\textbf{Legal Notice}
    \vspace{1em}

    \textbf{PROPRIETARY AND CONFIDENTIAL}
    \end{center}

    \noindent This document, including all algorithms, models, methodology, formulas,
    calculations, and intellectual concepts contained herein, constitutes the exclusive
    intellectual property and confidential trade secrets of MKM Research Labs (``MKM'').

    \vspace{0.5em}
    \noindent Any usage, reproduction, distribution, modification, or disclosure of this
    document or any portion thereof, without the express prior written authorisation from
    MKM Research Labs is strictly prohibited and constitutes an infringement of intellectual
    property rights.

    \vspace{0.5em}
    \noindent The document is provided on an ``as is'' basis. MKM Research Labs makes no
    representations or warranties of any kind, express or implied, regarding its accuracy,
    reliability, or suitability for any particular purpose.

    \vspace{1em}
    \begin{center}
    All rights reserved. \copyright\ 2021--2026 MKM Research Labs.

    \vspace{0.5em}
    Governed by the laws of the United Kingdom.
    \end{center}
    \clearpage
}

\newcommand{\mkmhistorypage}[1]{%
    \section{Document History}
    \begin{table}[ht]
        \centering
        \begin{tabular}{l|l|l|l}
            \textbf{Date} & \textbf{Version} & \textbf{Description} & \textbf{Author} \\
            \hline
            #1
        \end{tabular}
        \label{tab:revision_history}
    \end{table}
    \clearpage
}

% ─── Bank validation protocol section commands ──────────────────────────────
% Use these to ensure every model doc has the required sections.
% They are just styled \section calls — content is provided by each doc.

\newcommand{\mkmexecsummary}{\section{Executive Summary}}
\newcommand{\mkmpurpose}{\section{Model Purpose and Scope}}
\newcommand{\mkmframework}{\section{Mathematical Framework}}
\newcommand{\mkminputs}{\section{Input Parameters and Calibration}}
\newcommand{\mkmimplementation}{\section{Implementation Details}}
\newcommand{\mkmsensitivity}{\section{Sensitivity Analysis}}
\newcommand{\mkmvalidation}{\section{Validation and Backtesting}}
\newcommand{\mkmtestresults}{\subsection{Automated Test Results}\input{test_results}}
\newcommand{\mkmlimitations}{\section{Model Limitations and Known Weaknesses}}

\newcommand{\mkmgovernance}{\section{Governance Validation Questions}
\noindent The following nine questions are mandated by Chapter~5 of the
\emph{Model Risk Governance for Vendors} handbook and must be explicitly
addressed for every model in the inventory. Response status is tracked in the
Model Governance Dashboard (MRC portal). \par\medskip}

\newcommand{\mkmvqitem}[3]{\subsection{Q#1 --- #2}
\noindent\textit{``#3''}\par\medskip\noindent}

\newcommand{\mkmchangelog}{\section{Change History}}

% T1 so the underscores in case names and code identifiers print as underscores.
\usepackage[T1]{fontenc}
\usepackage{lmodern}
% Generated by tools.wolfram_oracle -- do not edit.
\newcommand{\bmGenerated}{2026-10-19}
\newcommand{\bmCodeFile}{\texttt{src/models/hazard/gev.py}}
\newcommand{\bmCodeSha}{\texttt{cca7d20506171a85}}
\newcommand{\bmCodeCommit}{\texttt{8bbb73347e54}}
\newcommand{\bmCodeDirty}{}
\newcommand{\bmNumerics}{numpy 1.26.4, scipy 1.15.1, mpmath 1.4.1}
\newcommand{\bmWolfram}{15.0.0 for Mac OS X ARM (64-bit) (May 26, 2026)}
\newcommand{\bmWolframVia}{ via live kernel session (wolframclient)}
\newcommand{\bmWolframConfirms}{Wolfram and the 50-digit reference agree to $5.0\times10^{-30}$, so the reference is confirmed by a second, independent implementation.}
\newcommand{\bmData}{synthetic samples from known parameters, seed 20261019}
\newcommand{\bmTol}{$1.0\times10^{-9}$}
\newcommand{\bmOperatingT}{10{,}000}
\newcommand{\bmWorstLevel}{$3.0\times10^{-11}$}
\newcommand{\bmWorstLevelAt}{below\_gumbel\_switch/T=10000}
\newcommand{\bmWorstExc}{$4.2\times10^{-10}$}
\newcommand{\bmWorstExcAt}{below\_gumbel\_switch/T=10000}
\newcommand{\bmLLGain}{$2.1\times10^{-6}$}
\newcommand{\bmLLTol}{$1.0\times10^{-4}$}
\newcommand{\bmRatioMin}{1.02}
\newcommand{\bmRatioMax}{4.46}
\newcommand{\bmTrapRight}{398.07}
\newcommand{\bmTrapWrong}{174.88}
\newcommand{\bmChecks}{6}
\newcommand{\bmOutcome}{all pass}


\begin{document}

\mkmtitlepage
\mkmlegalpage

\tableofcontents
\clearpage

% =========================================================================
\mkmexecsummary
% =========================================================================

This document benchmarks the implementation of the GEV Hazard Curve Model
(MKM-GH-001): the \texttt{GEVFitter} class that fits a generalised extreme value
(GEV) distribution to gauge peak levels and turns the fit into return levels and
exceedance probabilities. It asks two questions of the code, and answers each
against references that share none of its numerics.

\begin{enumerate}[leftmargin=*]
\item \textbf{Does it compute its formulae correctly?} Return levels and
  exceedance probabilities are checked against the closed form evaluated to 50
  significant digits, and against Wolfram. Within the operating range
  ($T \le \bmOperatingT$ years) the worst relative error is \bmWorstLevel{} on a
  return level and \bmWorstExc{} on an exceedance probability, against a
  tolerance of \bmTol. \bmWolframConfirms
\item \textbf{Does its fit reach the maximum of the likelihood?} The model's
  maximum-likelihood fit is compared with an independent multi-start optimiser
  on a likelihood written from the density. The largest gain in log-likelihood
  the independent fit finds is \bmLLGain{} (tolerance \bmLLTol): the model's
  optimiser stops at the maximum.
\end{enumerate}

\paragraph{Finding.} The implementation is sound; the tail risk is in the
\emph{estimate}, not the arithmetic. On samples drawn from known parameters, the
fitted 1-in-1000 level ranges from \bmRatioMin{} to \bmRatioMax{} times the true
one through sampling variation alone (Table~\ref{tab:bm-fit}). The model reports
return levels without a confidence interval.

\paragraph{Evidence.} \bmChecks{} checks, outcome: \bmOutcome{}
(Table~\ref{tab:bm-checks}).
\begin{itemize}[leftmargin=*,nosep]
\item Wolfram: \bmWolfram{}\bmWolframVia.
\item Code under test: \bmCodeFile{} at commit \bmCodeCommit{}\bmCodeDirty.
\item SHA-256 of that file: \bmCodeSha{}\ldots
\item Numerics: \bmNumerics.
\item Generated \bmGenerated{} by \texttt{python -m tools.wolfram\_oracle gh001}.
\end{itemize}

% =========================================================================
\mkmpurpose
% =========================================================================

\subsection{What is benchmarked}
The four public methods of \texttt{GEVFitter} in
\texttt{src/models/hazard/gev.py}, exactly as the hazard-curve builder calls
them: \texttt{fit} (GEV, and Gumbel when forced), \texttt{return\_level},
\texttt{exceedance\_probability}, and the sign conversion between the
extreme-value shape $\xi$ and scipy's $c=-\xi$. The class is loaded from its file
alone, so the benchmark exercises the production code and nothing around it.

\subsection{What is not}
This is a benchmark of the implementation (validation question Q8). It is not
independent validation (Q9): it shows the model computes what its formulae say,
not that a GEV fitted to per-storm peaks is the right model for the hazard.

% =========================================================================
\mkmframework
% =========================================================================

With location $\mu$, scale $\sigma>0$ and shape $\xi$ ($\xi>0$ a heavy Fr\'echet
tail, $\xi<0$ bounded above), the GEV distribution function is
\begin{equation}
F(x) = \exp\!\left\{-\left[1+\xi\,\frac{x-\mu}{\sigma}\right]^{-1/\xi}\right\},
\qquad
F(x) = \exp\!\left\{-e^{-(x-\mu)/\sigma}\right\}\ \ (\xi=0).
\end{equation}
The $T$-year return level solves $F(x_T)=1-1/T$. Writing the reduced variate
$y_T=-\log(1-1/T)$,
\begin{equation}
x_T = \mu + \frac{\sigma}{\xi}\left(y_T^{-\xi}-1\right),
\qquad x_T = \mu - \sigma\log y_T\ \ (\xi=0),
\end{equation}
and the exceedance probability is $P(X>x) = 1-F(x)$. \texttt{GEVFitter} takes the
Gumbel branch when $|\xi|<10^{-10}$.

The log-likelihood of peaks $x_1,\dots,x_n$, with $t_i = 1+\xi(x_i-\mu)/\sigma>0$, is
\begin{equation}
\ell(\mu,\sigma,\xi) = -n\log\sigma - \left(1+\frac1\xi\right)\sum_i \log t_i
 - \sum_i t_i^{-1/\xi}.
\end{equation}

\paragraph{Sign convention.} scipy's \texttt{genextreme} uses $c=-\xi$. Passed
unflipped, the 1-in-1000 level for $\xi=0.2$ would be \bmTrapWrong{} rather than
\bmTrapRight; \texttt{GEVFitter} converts in one named place and is correct.

% =========================================================================
\section{Benchmark Method}
% =========================================================================

\subsection{Closed-form reference}
Each return level and exceedance probability is evaluated with \texttt{mpmath}
at 50 significant digits, using $\mathrm{expm1}$ and the reduced variate so
neither $1-1/T$ nor $1-F$ loses digits. The exceedance threshold for each case
is the exact $T$-year level, so the right answer is $1/T$. Six parameter cases
cover the Gumbel, Fr\'echet, bounded and near-zero-shape regimes, including a
shape just below the model's Gumbel switch.

\subsection{Wolfram oracle}
\begin{sloppypar}
The same quantities are evaluated in Wolfram with \texttt{Quantile} and
\texttt{SurvivalFunction} of \texttt{MaxStableDistribution}$[\mu,\sigma,\xi]$
(same sign of $\xi$) and, for $\xi=0$, \texttt{ExtremeValueDistribution}.
Wolfram's \texttt{GumbelDistribution} is the \emph{minimum} law and would be
wrong here.
Inputs go to Wolfram as the exact value of each double the model used, set to
60-digit precision, so \texttt{N[$\cdot$, 30]} returns 30 genuine digits rather
than machine precision.
\end{sloppypar}

\subsection{Independent fit}
The likelihood above is maximised by Nelder--Mead from six starting shapes, with
no \texttt{scipy.stats} code, and evaluated at 50 digits at both the model's and
the independent parameters. Wolfram's \texttt{EstimatedDistribution} gives a
third fit. A positive gain means the model's optimiser stopped short of the
maximum. Data: \bmData.

% =========================================================================
\mkmvalidation
% =========================================================================

% Generated by tools.wolfram_oracle -- do not edit.
\begin{table}[H]\centering\small
\caption{Benchmark checks}\label{tab:bm-checks}
\begin{adjustbox}{max width=\textwidth}
\begin{tabular}{lrrl}
\toprule
Check & Result & Tolerance & Outcome \\
\midrule
Return levels{,} $T \le 10{,}000$ & $3.0\times10^{-11}$ & $1.0\times10^{-9}$ & pass \\
Exceedance probabilities{,} $T \le 10{,}000$ & $4.2\times10^{-10}$ & $1.0\times10^{-9}$ & pass \\
Wolfram agrees with the reference & $5.0\times10^{-30}$ & $1.0\times10^{-9}$ & pass \\
Model agrees with Wolfram & $4.2\times10^{-10}$ & $1.0\times10^{-9}$ & pass \\
Fit at the maximum: independent optimiser & $2.1\times10^{-6}$ & $1.0\times10^{-4}$ & pass \\
Fit at the maximum: Wolfram & $2.1\times10^{-6}$ & $1.0\times10^{-4}$ & pass \\
\bottomrule
\end{tabular}
\end{adjustbox}
\end{table}


\subsection{Return levels and exceedance}
Worst in range: return level at \texttt{\bmWorstLevelAt}; \\
exceedance at \texttt{\bmWorstExcAt}. The \texttt{below\_gumbel\_switch} case carries the
model's rounding to Gumbel at $\xi=10^{-11}$, an error of order $\xi$ that is
immaterial. Beyond the
operating range the model's $\mathrm{ppf}(1-p)$ and $1-\mathrm{cdf}$ lose
digits as $p\to0$; this is reported, not judged.

% Generated by tools.wolfram_oracle -- do not edit.
{\footnotesize\setlength{\tabcolsep}{4pt}\begin{longtable}{lrrrrrr}
\caption{GEVFitter against the 50-digit reference: relative errors. The exceedance threshold is the exact $T$-year level, so $P(X>x_T)=1/T$}\label{tab:bm-return-levels}\\
\toprule
Case & $T$ & Reference level & Level err & Exc.\ err & Level vs W. & Exc.\ vs W. \\
\midrule
\endhead
gumbel & 10 & 145.007346546248 & $1.3\times10^{-16}$ & $6.4\times10^{-16}$ & $1.7\times10^{-16}$ & $6.2\times10^{-16}$ \\
gumbel & 100 & 192.002984535531 & $2.3\times10^{-17}$ & $6.7\times10^{-16}$ & $2.4\times10^{-19}$ & $6.8\times10^{-16}$ \\
gumbel & 200 & 205.916242850700 & $2.1\times10^{-17}$ & $6.8\times10^{-16}$ & $8.4\times10^{-17}$ & $6.7\times10^{-16}$ \\
gumbel & 1{,}000 & 238.145101410474 & $2.0\times10^{-17}$ & $6.4\times10^{-16}$ & $6.9\times10^{-21}$ & $6.6\times10^{-16}$ \\
gumbel & 10{,}000 & 284.205807397854 & $7.7\times10^{-15}$ & $1.1\times10^{-13}$ & $7.8\times10^{-15}$ & $1.1\times10^{-13}$ \\
gumbel & 1{,}000{,}000 & 376.310201159281 & $1.5\times10^{-12}$ & $2.9\times10^{-11}$ & $1.5\times10^{-12}$ & $2.9\times10^{-11}$ \\
frechet\_moderate & 10 & 156.842740650253 & $1.2\times10^{-17}$ & $1.6\times10^{-16}$ & $4.0\times10^{-17}$ & $1.4\times10^{-16}$ \\
frechet\_moderate & 100 & 250.936528171715 & $2.0\times10^{-16}$ & $1.0\times10^{-15}$ & $2.1\times10^{-16}$ & $1.0\times10^{-15}$ \\
frechet\_moderate & 200 & 288.395445935869 & $3.9\times10^{-17}$ & $6.9\times10^{-16}$ & $4.3\times10^{-17}$ & $6.8\times10^{-16}$ \\
frechet\_moderate & 1{,}000 & 398.067345230813 & $7.2\times10^{-17}$ & $5.4\times10^{-16}$ & $1.0\times10^{-16}$ & $5.5\times10^{-16}$ \\
frechet\_moderate & 10{,}000 & 630.951034675384 & $2.2\times10^{-14}$ & $1.1\times10^{-13}$ & $2.2\times10^{-14}$ & $1.1\times10^{-13}$ \\
frechet\_moderate & 1{,}000{,}000 & 1584.89303397173 & $5.8\times10^{-12}$ & $2.9\times10^{-11}$ & $5.8\times10^{-12}$ & $2.9\times10^{-11}$ \\
frechet\_heavy & 10 & 183.231304990444 & $3.3\times10^{-17}$ & $1.3\times10^{-16}$ & $1.6\times10^{-17}$ & $1.1\times10^{-16}$ \\
frechet\_heavy & 100 & 458.997067605108 & $5.0\times10^{-16}$ & $8.8\times10^{-16}$ & $5.3\times10^{-16}$ & $9.0\times10^{-16}$ \\
frechet\_heavy & 200 & 624.977284199435 & $2.8\times10^{-16}$ & $1.1\times10^{-15}$ & $3.5\times10^{-16}$ & $1.1\times10^{-15}$ \\
frechet\_heavy & 1{,}000 & 1324.59474401879 & $2.6\times10^{-16}$ & $1.1\times10^{-15}$ & $3.4\times10^{-16}$ & $1.1\times10^{-15}$ \\
frechet\_heavy & 10{,}000 & 4059.89999708317 & $5.5\times10^{-14}$ & $1.1\times10^{-13}$ & $5.5\times10^{-14}$ & $1.1\times10^{-13}$ \\
frechet\_heavy & 1{,}000{,}000 & 40059.9899999970 & $1.4\times10^{-11}$ & $2.9\times10^{-11}$ & $1.4\times10^{-11}$ & $2.9\times10^{-11}$ \\
weibull\_bounded & 10 & 136.241869030462 & $8.7\times10^{-17}$ & $4.5\times10^{-16}$ & $9.0\times10^{-17}$ & $4.2\times10^{-16}$ \\
weibull\_bounded & 100 & 160.149285268037 & $2.6\times10^{-17}$ & $3.6\times10^{-16}$ & $4.4\times10^{-17}$ & $3.7\times10^{-16}$ \\
weibull\_bounded & 200 & 165.325388660181 & $4.4\times10^{-17}$ & $1.7\times10^{-16}$ & $1.9\times10^{-17}$ & $1.8\times10^{-16}$ \\
weibull\_bounded & 1{,}000 & 174.878622625521 & $5.5\times10^{-18}$ & $1.1\times10^{-15}$ & $4.8\times10^{-18}$ & $1.1\times10^{-15}$ \\
weibull\_bounded & 10{,}000 & 184.150909578672 & $2.0\times10^{-15}$ & $1.1\times10^{-13}$ & $2.0\times10^{-15}$ & $1.1\times10^{-13}$ \\
weibull\_bounded & 1{,}000{,}000 & 193.690425924240 & $1.9\times10^{-13}$ & $2.9\times10^{-11}$ & $1.9\times10^{-13}$ & $2.9\times10^{-11}$ \\
near\_zero\_shape & 10 & 145.007346596890 & $4.2\times10^{-17}$ & $6.6\times10^{-17}$ & $2.2\times10^{-17}$ & $8.8\times10^{-17}$ \\
near\_zero\_shape & 100 & 192.002984747145 & $5.1\times10^{-17}$ & $4.0\times10^{-16}$ & $4.9\times10^{-17}$ & $4.1\times10^{-16}$ \\
near\_zero\_shape & 200 & 205.916243131156 & $9.1\times10^{-17}$ & $1.4\times10^{-15}$ & $9.9\times10^{-17}$ & $1.4\times10^{-15}$ \\
near\_zero\_shape & 1{,}000 & 238.145101887576 & $1.6\times10^{-17}$ & $1.1\times10^{-15}$ & $1.2\times10^{-17}$ & $1.1\times10^{-15}$ \\
near\_zero\_shape & 10{,}000 & 284.205808246148 & $7.6\times10^{-15}$ & $1.1\times10^{-13}$ & $7.6\times10^{-15}$ & $1.1\times10^{-13}$ \\
near\_zero\_shape & 1{,}000{,}000 & 376.310203067964 & $1.5\times10^{-12}$ & $2.9\times10^{-11}$ & $1.5\times10^{-12}$ & $2.9\times10^{-11}$ \\
below\_gumbel\_switch & 10 & 145.007346546755 & $3.5\times10^{-12}$ & $2.4\times10^{-11}$ & $3.5\times10^{-12}$ & $2.4\times10^{-11}$ \\
below\_gumbel\_switch & 100 & 192.002984537647 & $1.1\times10^{-11}$ & $1.1\times10^{-10}$ & $1.1\times10^{-11}$ & $1.1\times10^{-10}$ \\
below\_gumbel\_switch & 200 & 205.916242853505 & $1.4\times10^{-11}$ & $1.4\times10^{-10}$ & $1.4\times10^{-11}$ & $1.4\times10^{-10}$ \\
below\_gumbel\_switch & 1{,}000 & 238.145101415245 & $2.0\times10^{-11}$ & $2.4\times10^{-10}$ & $2.0\times10^{-11}$ & $2.4\times10^{-10}$ \\
below\_gumbel\_switch & 10{,}000 & 284.205807406337 & $3.0\times10^{-11}$ & $4.2\times10^{-10}$ & $3.0\times10^{-11}$ & $4.2\times10^{-10}$ \\
below\_gumbel\_switch & 1{,}000{,}000 & 376.310201178368 & $5.2\times10^{-11}$ & $9.7\times10^{-10}$ & $5.2\times10^{-11}$ & $9.7\times10^{-10}$ \\
\bottomrule
\end{longtable}}


\subsection{Maximum-likelihood fit}
% Generated by tools.wolfram_oracle -- do not edit.
\begin{table}[H]\centering\small
\caption{Maximum-likelihood fit: GEVFitter against an independent optimiser}\label{tab:bm-fit}
\begin{adjustbox}{max width=\textwidth}
\begin{tabular}{llrrrrrrrrr}
\toprule
Series & Fit & $n$ & $\hat\xi$ model & $\hat\xi$ indep. & $\xi$ true & LL gain & $x_{1000}$ model & $x_{1000}$ true & $\hat\xi$ Wolfram & LL gain W. \\
\midrule
n30\_xi0.1 & gev & 30 & 0.2449 & 0.2449 & 0.10 & $1.1\times10^{-7}$ & 9.741 & 6.981 & 0.2449 & $1.1\times10^{-7}$ \\
n30\_xi0.1 & gumbel & 30 & 0.0000 & 0.0000 & 0.10 & $-4.7\times10^{-15}$ & 6.011 & 6.981 & 0.0000 & $-2.0\times10^{-14}$ \\
n50\_xi0.3 & gev & 50 & 0.6588 & 0.6588 & 0.30 & $9.0\times10^{-8}$ & 54.618 & 12.256 & 0.6588 & $9.0\times10^{-8}$ \\
n50\_xi0.3 & gumbel & 50 & 0.0000 & 0.0000 & 0.30 & $-1.8\times10^{-15}$ & 7.262 & 12.256 & 0.0000 & $4.2\times10^{-14}$ \\
n200\_xi-0.1 & gev & 200 & -0.0832 & -0.0833 & -0.10 & $2.1\times10^{-6}$ & 5.087 & 4.995 & -0.0833 & $2.1\times10^{-6}$ \\
n200\_xi-0.1 & gumbel & 200 & 0.0000 & 0.0000 & -0.10 & $-5.9\times10^{-15}$ & 5.659 & 4.995 & 0.0000 & $1.7\times10^{-14}$ \\
n30\_xi0.4 & gev & 30 & 0.7073 & 0.7073 & 0.40 & $1.0\times10^{-7}$ & 76.048 & 17.846 & 0.7073 & $1.0\times10^{-7}$ \\
n30\_xi0.4 & gumbel & 30 & 0.0000 & 0.0000 & 0.40 & $-2.9\times10^{-15}$ & 8.996 & 17.846 & 0.0000 & $-2.1\times10^{-14}$ \\
\bottomrule
\end{tabular}
\end{adjustbox}
\end{table}


% =========================================================================
\mkmlimitations
% =========================================================================

\begin{enumerate}[leftmargin=*]
\item \textbf{Benchmark, not validation.} The model is shown to compute its
  formulae; whether those formulae suit the hazard is a separate question.
\item \textbf{Per-storm peaks.} \texttt{GEVFitter} fits the GEV to every
  storm's peak level. The GEV is the limiting law for block maxima (e.g.\ annual
  maxima); the per-storm choice is a methodology question this benchmark does
  not answer.
\item \textbf{Parameter uncertainty dominates.} At the sample sizes tested,
  sampling variation moves the 1-in-1000 level by a factor of up to \bmRatioMax;
  the model reports no confidence interval.
\item \textbf{Data.} Unless a gauge export is named under Evidence, the fit
  section uses synthetic samples; it then shows estimator behaviour, not the
  behaviour on the real gauge record.
\end{enumerate}

\mkmhistorypage{19-Oct-2026 & 0.1 & First issue & David K Kelly \\}

\end{document}
